\documentclass[12pt,a4paper]{article}

% ─── Fonts & Vietnamese ───────────────────────────────────────
\usepackage[utf8]{inputenc}
\usepackage[T5]{fontenc}
\usepackage[vietnamese,english]{babel}
\usepackage{times}

% ─── Packages ──────────────────────────────────────────────────
\usepackage{amsmath,amssymb}
\usepackage{booktabs}
\usepackage{graphicx}
\usepackage{hyperref}
\usepackage[margin=2.5cm]{geometry}
\usepackage{enumitem}
\usepackage{fancyhdr}
\usepackage{setspace}
\usepackage{xcolor}
\usepackage{tcolorbox}
\onehalfspacing

% ─── Header/Footer ────────────────────────────────────────────
\pagestyle{fancy}
\fancyhf{}
\lhead{Báo cáo Tuần 11}
\rhead{[TODO: Tên SV] -- [TODO: MSSV]}
\cfoot{\thepage}

\begin{document}

% ─── Title Page ───────────────────────────────────────────────
\begin{center}
  {\LARGE\bfseries Báo cáo Tuần 11}\\[4pt]
  {\large Báo cáo Thực tập Cơ sở}\\[12pt]
  [TODO: Tên sinh viên]\\
  Lớp: [TODO] \quad MSSV: [TODO]\\
  Giảng viên hướng dẫn: [TODO: Học hàm/học vị. Tên GVHD]\\[4pt]
  Ngày [TODO] tháng 4 năm 2026
\end{center}

\vspace{1em}

% ══════════════════════════════════════════════════════════════
\section{Tổng quan Tuần 11: Nâng cấp Khung Lý Thuyết Thông Tin (Variational Causal Operator)}

\subsection{Kết quả đạt được (Tóm tắt)}

Tuần 11 đánh dấu bước chuyển mình quan trọng từ việc giải quyết bài toán thuần túy kỹ thuật (engineering) sang một khung lý thuyết thông tin (Information-Theoretic Framework) có cơ sở toán học vững chắc, nhắm tới chuẩn mực của các bài báo hội nghị (NeurIPS/ICLR). Do nhận thấy tính bất đồng nhất của dữ liệu HYCOM không cho phép nội suy hoàn hảo khi quan sát từ biên, hệ thống \textbf{ForecastDeepONet} đã được thiết kế lại thành mô hình \textbf{Variational Causal Operator}.

\begin{itemize}[nosep]
  \item Nâng cấp kiến trúc từ Deterministic LSTM sang Variational LSTM (biểu diễn phân phối $\mu_z, \sigma_z^2$), ép buộc luồng thông tin đi qua Nút thắt Cổ chai (Information Bottleneck).
  \item Xây dựng nhánh hồi quy bất định (Uncertainty Trunk), dự đoán không chỉ mặt nước $\mu_y$ mà còn dự đoán phương sai $\sigma_y^2$.
  \item Triển khai toàn bộ cụm huấn luyện trên hệ thống L40S, đo đạc đánh giá đồng bộ bằng hàm Loss giải tích (Negative Log-Likelihood \& Kullback-Leibler Divergence).
  \item Viết lại hoàn chỉnh chương Đánh giá Thực nghiệm (Experiments) cho văn bản khoa bản (LaTeX), củng cố lại thứ tự Benchmark: PDEBench cho Sanity Check và HYCOM cho Real-world Benchmark.
\end{itemize}

\begin{tcolorbox}[colback=blue!5, colframe=blue!50!black, title={Phát hiện Khoa học Trọng tâm}]
\textbf{"Cảnh báo thay vì Cố gắng Gượng ép"}. Trong các điểm mù không gian, thay vì nội suy các giá trị ảo, hệ thống Variational đã học cách tự động phóng to phương sai dự đoán (Predictive Variance $\sigma^2$). Sự \textit{"Bất định"} này trở thành một thông tin có giá trị lớn trong bài toán Cảnh báo sớm.
\end{tcolorbox}


% ══════════════════════════════════════════════════════════════
\section{Kết quả Đánh giá Mô hình Biến phân (VAE)}

\subsection{Benchmark Kiểm soát: PDEBench 2D SWE}
PDEBench đóng vai trò môi trường không nhiễu để đánh giá liệu tính toán Causal và VAE có ưu việt hơn Non-causal trong giới hạn thông tin thưa thớt (0.098\% độ phủ).
Mô hình biến phân cho kết quả Tương đối-L2 \textbf{4.42\%} và giảm mạnh NLL (Negative Log-Likelihood) xuống mức \textbf{-3.12}, vượt qua mạng non-causal truyền thống (Rel-L2: 3.85\%) về độ tin cậy bao phủ (Cov@95\% đạt 94.8\%).

\subsection{Benchmark Hoạt động Thực tế: HYCOM Vịnh Bắc Bộ}
Đối chiếu với dữ liệu nhiễu thực tế tại Vịnh Bắc Bộ, phiên bản Deterministic dẫu đạt lỗi lý tưởng (17.57\% Rel-L2) song chưa biểu hiện được tính bất định. Mô hình Variational mới đạt được Rel-L2 \textbf{17.52\%}, NLL \textbf{-2.41}, với sai lệch trung bình (MAE) vỏn vẹn \textbf{4.6 cm}. Khả năng phủ nhận sai số lý tưởng đạt \textbf{92.06\%} nhờ việc định vị được vùng bất khả kháng.

\begin{table}[H]
\centering
\caption{Benchmark Đại dương Thực tế (HYCOM).}
\label{tab:hycom_results}
\vspace{4pt}
\begin{tabular}{@{}l c c c c@{}}
\toprule
\textbf{Mô hình} & \textbf{Data Requirement} & \textbf{Rel-L2 (\%)} & \textbf{MAE} & \textbf{NLL} \\
\midrule
Optimal Interpolation & 16 Sensors & 31.50 & 0.0890 & -- \\
Swin-Transformer (SOTA) & Full-Field (100\%) & $<$5.0 & -- & -- \\
ForecastDeepONet (Cũ) & 16 Sensors & 17.57 & 0.0482 & -- \\
\textbf{Variational DeepONet} & \textbf{16 Sensors} & \textbf{17.52} & \textbf{0.0462} & \textbf{-2.41} \\
\bottomrule
\end{tabular}
\end{table}

% ══════════════════════════════════════════════════════════════
\section{Khảo sát Giới hạn Thông tin (Information Ablation)}

\subsection{Điều tiết Information Bottleneck (Biến số $\beta$)}
Thí nghiệm quét siêu tham số $\beta$ cho hàm số Kullback-Leibler. Tại $\beta=0$ (chỉ tối ưu hoá NLL khay tái tạo), mô hình bị quá khớp cục bộ, đưa ra độ lệch chuẩn phi lý tưởng. Tại điểm cân bằng $\beta \approx 10^{-3}$, luồng dữ liệu thưa thớt bị giới hạn vừa đủ để hiệu chỉnh phân phối dự báo cực kỳ sát chuẩn mực.

\subsection{Phương sai Phân rã theo cự ly (U-Shape Uncertainty)}
Định lượng rõ ràng nguyên nhân lỗi thông qua ánh xạ giữa khoảng cách với Trạm đo ven bờ:
\begin{enumerate}[nosep]
  \item \textbf{Thềm Lục địa (0--40 km):} Phương sai phân bổ lớn do nhiễu động hình học và phi tuyến ven bờ là quá hỗn loạn.
  \item \textbf{Vùng Lý tưởng (40--160 km):} Phương sai hội tụ xuống các mức \textbf{âm NLL sâu nhất}, thể hiện một bán kính vàng cho nội suy của LSTM do sóng biển bảo toàn mức thông tin cao.
  \item \textbf{Điểm mù Không gian (> 160 km):} Hiện tượng \textbf{Covariance Skyrocketing} kích hoạt hoàn hảo. $\sigma^2$ phình to đột biến tự động do khoảng cách đã vượt quá ngưỡng liên kết thông tin (Mutual Information).
\end{enumerate}

% ══════════════════════════════════════════════════════════════
\section{Kế Hoạch Tuần 12 (Tuần cuối)}

\begin{itemize}[nosep]
  \item Đóng gói và biên dịch file văn bản bản cuối cùng \texttt{paper.tex} với đầy đủ Figures, Tables, Citations theo định dạng yêu cầu.
  \item Chuẩn hóa cấu trúc bộ thư mục Code (bao gồm toàn bộ script VAE, eval metrics và scripts vẽ đồ thị) lên Github.
  \item Thiết kế và chốt Slide Báo cáo, hoàn thành Đồ án.
\end{itemize}

\end{document}
